\appendix

\section{Per-Task Accuracy Progression}
\label{app:pertask}

Table~\ref{tab:pertask} shows task-by-task accuracy for \methodname{}+SHINE and the SHINE baseline.

\begin{table}[h]
\centering
\caption{Per-task accuracy progression on the default Marathon CIL ordering.}
\label{tab:pertask}
\begin{tabular}{clcccc}
\toprule
Task & Domain & \#Seen & AnchorLoRA+S & SHINE & $\Delta$ \\
\midrule
0 & Trento & 2 & 100.0 & 100.0 & 0.0 \\
1 & Trento & 4 & 97.9 & 97.9 & 0.0 \\
2 & Trento & 6 & 94.1 & 94.1 & 0.0 \\
3 & Houston & 11 & 94.2 & 91.5 & +2.7 \\
4 & Houston & 16 & 88.5 & 82.7 & +5.8 \\
5 & Houston & 21 & 84.9 & 74.3 & +10.6 \\
6 & MUUFL & 25 & 70.8 & 66.9 & +3.9 \\
7 & MUUFL & 29 & 76.3 & 68.3 & +8.0 \\
8 & MUUFL & 32 & 74.4 & 66.1 & +8.3 \\
\bottomrule
\end{tabular}
\end{table}

\section{Forgetting Analysis}
\label{app:forgetting}

\begin{table}[h]
\centering
\caption{Per-domain forgetting for \methodname{}+SHINE.}
\label{tab:forgetting}
\begin{tabular}{lccc}
\toprule
Domain & Peak Acc & Final Acc & Forgetting \\
\midrule
Trento & 100.0 (T0) & 91.7 (T8) & $-$8.3pp \\
Houston & 92.5 (T3) & 67.9 (T8) & $-$24.6pp \\
MUUFL & 79.7 (T7) & 74.0 (T8) & $-$5.7pp \\
\bottomrule
\end{tabular}
\end{table}

Houston exhibits the largest forgetting ($-$24.6pp), primarily at Task 6 when the MUUFL domain enters. This cross-domain forgetting pattern is observed across all methods and represents a fundamental challenge of cross-domain marathon CIL rather than a method-specific limitation.

\section{Negative Results Summary}
\label{app:negative}

We tested 12 architectural variants. The following designs were eliminated:
\begin{itemize}
    \item \textbf{DualSpeed-S$^2$CM}: Dual classification heads (spectral NCM + fusion NCM) with calibrated blending. Result: 72.3\% ($-$3.1pp). The spectral-only head is too weak (64-d features insufficient for 32-class discrimination).
    \item \textbf{SpecRoute}: Spectral-routed expert banks where frozen spectral features select spatial LoRA experts via soft attention. Result: 72.1\% ($-$3.3pp). Spectral centroids lack discriminability across tasks within the same domain.
    \item \textbf{Orthogonal LoRA}: InfLoRA-style null-space expansion. Result: 74.4\% ($-$1.0pp). Rank-4 subspace is too constrained for orthogonal partitioning.
    \item \textbf{Auto-freeze}: Drift-based automatic spectral freeze detection. Spectral drift never converges during training because new classes continuously modify the representation. Fixed warmup is more practical.
\end{itemize}
